\chapter{Implementation}
\label{chap:implementation}

\section{Microcontroller Firmware Architecture}
The edge sensing firmware is developed in C/C++ on top of the Espressif ESP-IDF framework within the Arduino development ecosystem. The firmware initializes the Wi-Fi physical layer in Station (STA) mode and registers the native CSI hardware interrupt hook via \texttt{esp\_wifi\_set\_csi\_rx\_cb()}. 

Upon receipt of each valid 802.11n OFDM packet, the baseband hardware invokes this callback function, populating a \texttt{wifi\_csi\_info\_t} structure containing the packet RSSI, noise floor, timestamp, and 52 signed In-phase ($I$) and Quadrature ($Q$) subcarrier bytes. To maintain a deterministic, continuous packet flow without relying on irregular user network browsing, the ESP32 transmits lightweight UDP echo ping packets to the local gateway at a calibrated frequency of 20~Hz. If the Wi-Fi connection is temporarily interrupted, the firmware automatically falls back into promiscuous sniffing mode, ensuring uninterrupted CSI acquisition from ambient broadcast packets.

\section{Non-Blocking USB Serial Streaming Pipeline}
Each raw CSI frame is formatted into a compact serial packet:
\begin{verbatim}
CSI_DATA,<frame_id>,<rssi>,<I_0>,<Q_0>,<I_1>,<Q_1>,...,<I_51>,<Q_51>
\end{verbatim}
The packet is dispatched across the CP2102/CH340 USB-UART bridge at a baud rate of 115200 baud. On the host receiving machine, a non-blocking asynchronous reader thread buffers arriving serial lines, performs CRC integrity checks, and passes complete 52-subcarrier arrays to the DSP pipeline via thread-safe FIFO queues, eliminating serial read latency.

\section{Real-Time Python DSP Processing Engine}
The core DSP processing engine is implemented in Python utilizing high-performance vectorized operations with NumPy and SciPy. The pipeline executes the following sequential steps for each incoming frame:
\begin{enumerate}
    \item \textbf{Complex IQ Reconstitution:} Unpacks 52 pairs of integer $(I_k, Q_k)$ into complex amplitudes $|H_k| = \sqrt{I_k^2 + Q_k^2}$ and phases $\theta_k = \operatorname{atan2}(Q_k, I_k)$.
    \item \textbf{Top-$K$ Variance Ranking:} Calculates temporal variance across a sliding buffer ($W=50$ frames), sorts subcarriers by variance, and selects the $K=10$ highest-energy indices.
    \item \textbf{Welford Online Accumulation:} Updates running mean $\mu_n$, sum of squares $S_n$, and running variance $\sigma_n^2$ in $O(1)$ memory.
    \item \textbf{Respiration Filtering and FFT:} Unwraps phase across time, applies the 2nd-order Butterworth bandpass filter (0.1--0.5~Hz), and computes FFT magnitude across a 20-second rolling window to locate the dominant respiratory peak.
    \item \textbf{Decision State Classification:} Evaluates the motion energy $V(t)$ against dynamic threshold $\tau_{th}$, outputting Room Empty, Subtle Motion, or Active Walking, alongside the 0--99\% confidence score.
\end{enumerate}

\section{Interactive Web Telemetry Dashboard}
To enable seamless real-time visualization and user monitoring, the host engine hosts a lightweight local HTTP server (port 8888). The frontend interface is constructed using HTML5, CSS3, and Chart.js. It features:
\begin{itemize}
    \item A dynamic 52-subcarrier amplitude bar graph displaying instantaneous channel frequency responses in real time.
    \item A rolling Welford variance timeline plotted alongside the auto-calibrated threshold.
    \item A live respiration waveform display accompanied by instantaneous BPM readouts.
    \item An activity status badge displaying presence state, confidence percentage, and connection health.
\end{itemize}
The dashboard updates seamlessly at 20~Hz with an HTTP telemetry latency of only 1.40~ms.
